\begin{lemma}[Stable change, external first-missing case]
Fix a finite loopless multigraph with its exact edge set, a palette of
$k$ colors, and a distinguished edge $e_0$.  Let $(T,\varphi_n)$ be a
positive-rung extended Tashkinov tree.  Let $\mathcal B$ be a hierarchy
base and let
\[
 T_n=T_{n,0}\subset T_{n,1}\subset\cdots\subset
 T_{n,q}\subset T_{n,q+1}=T
\]
be a good hierarchy $H$ under $\varphi_n$.  Write
\[
 A=T_{n,q}^{*}=\noderef{HierarchyStarredStage}(H,q),\qquad
 D_{n,q}=\noderef{HierarchyStoredColorsAt}(H,q).
\]
Assume that $T$ has maximum property, that $V(T)$ is not elementary,
that $\varphi_n$ is proper and leaves exactly $e_0$ uncolored, and that
these data form a \noderef{HierarchyCounterexample} $C$.

Assume, uniformly for every positive-rung maximum ETT on the same
graph, palette, and distinguished edge, with arbitrary certified input
and generating colorings and with rung number smaller than that of $T$,
the six lower-rung assertions of the technical theorem: elementarity,
terminal interchangeability, the unconditional PE stability-before-support
assertion, the predecessor-ladder stability assertion conditional on
first-appearance index greater than one, and the unconditional stored-pair
disjointness assertion, the terminal PE component assertion, the defective-endpoint
order assertion, and the mod-coloring assertion.  Assume also that every
same-rung maximum ETT with a good hierarchy having fewer than $q$
dividers is elementary.  Finally assume the large-color estimate
\[
                 2n+11\leq |\overline\varphi_n(T_n)|.          \tag{L}
\]

Put $\ell=|T|-|A|$, the number of exact tail additions above $A$, and
\[
 T(y_{\ell-1})=T\big|_{\min\{|A|+(\ell-1),|T|\}}.
\]
Assume (6.6), namely that $V(T(y_{\ell-1}))$ is elementary under
$\varphi_n$, and assume that $C$ is lexicographically minimal among all
hierarchy counterexamples on the fixed graph, palette, and edge whose
ETTs have the same rung number as $T$ and whose hierarchies have exactly
$q$ dividers, in the sense of
\noderef{HierarchyCounterexampleLess} and
\noderef{HierarchyCounterexampleMeasure}.

Let $\alpha\ne\beta$ be colors missing on $T(y_{\ell-1})$.  For a color
$c$, put
\[
 i_c=\noderef{HierarchyFirstMissingPosition}(\varphi_n,T,c),
 \qquad v_c=T.\mathit{vertex}(i_c).
\]
Let $Q$ be the full $(\alpha,\beta)$-component through an arbitrary
incidence-graph node $z$, and put
\[
 \sigma_n=\noderef{TwoColorComponentSwap}
       (G,\varphi_n,\alpha,\beta,z).
\]
Assume that every vertex of $A$ precedes $v_\alpha$, that
$v_\alpha\prec v_\beta$, that the colors $\alpha,\beta$ do not belong
to $D_{n,q}$, and that
\[
 \alpha\notin
 \{c:\exists t\in\operatorname{Fin}|T|,
        i_\alpha\leq t<i_\beta,
        \ \varphi_n(T.\mathit{edge}(t))=c\}.                 \tag{U}
\]
Thus (U) is the exact edge-identity assertion
$\alpha\notin\varphi_n\langle T(v_\beta)-T(v_\alpha)\rangle$.
The component $Q$ is arbitrary and may be the component through
$v_\alpha$ itself.  These are exactly the hypotheses in case 3 of the
paper's stable-change lemma.

Then all of the following hold.
\begin{enumerate}
\item The coloring $\sigma_n$ is
$(T_{n,q}^{*},D_n,\varphi_n)$-weakly stable in the exact sense of
\noderef{HierarchyWeaklyStable}.
\item The coloring $\sigma_n$ is a
$\varphi_n\bmod T_n$ coloring in the sense of \noderef{ModColoring}.
Either $A=T_n$, the non-strict endpoint which the tablet's positive-rung
ETT type intentionally does not repackage, or $A$ has an actual maximum
ETT presentation under $\sigma_n$.  In particular the second alternative
always holds when $q\geq1$.
\item If $q\geq1$, the same ordered base core and closure and the
stages
\[
 T_n=T_{n,0}\subset\cdots\subset T_{n,q}=A
\]
form a good hierarchy of an actual maximum ETT on $A$ under
$\sigma_n$.  Its divider count is $q-1$; at every retained index the
sets $D_{n,j}$ and the pairs $\Gamma^j_\eta$ agree exactly with the
original ones.  In addition,
\[
 T_{n,q}\text{ is }
 \left(\bigcup_{\eta\in D_{n,q}}\Gamma^{q-1}_\eta\right)^{-}
 \text{-closed under }\sigma_n.                              \tag{M}
\]
\end{enumerate}
The full tree $T$ also has an actual corresponding maximum ETT
presentation under $\sigma_n$, and the displayed hierarchy remains good
with the same ordered base, stages, dividers, $D$-sets, and all
paper-defined $\Gamma$-sets, in the precise sense of
\noderef{SameGoodHierarchy}.
\end{lemma}
\begin{proof}
Unpack $C$.  Its fields provide positivity, maximum property, the
displayed \noderef{TashkinovHierarchy}, its
\noderef{GoodHierarchy} witness, and non-elementarity.  All component
equalities and disjointness statements below are in
\noderef{TwoColorComponent}; consequently every edge node is an
individual element of the multigraph edge type, even when two edges
have the same endpoints.

We first record two facts used below.  The hypothesis (6.6)
and the initial-segment fields of the hierarchy make every displayed
stage through $A$ elementary under $\varphi_n$.  The core-inside-first, nested, and
terminal fields of \noderef{TashkinovHierarchy} compose to an ordered
initial-segment relation $A\preccurlyeq T$, so $|A|\leq|T|$.  If
$\ell=0$, the definition $\ell=|T|-|A|$ forces $|A|=|T|$; equal length
together with the ordered initial-segment
equalities identifies $A=T$ entry by entry.  At $\ell=0$ the prefix length
in (6.6) is
$\min\{|A|+(\ell-1),|T|\}=|T|$, because subtraction is truncated.
The full-length \noderef{TreeSequencePrefix} is $T$, so (6.6) would make
$V(T)$ elementary, contrary to the `notElementary' field of $C$.
Hence $\ell\geq1$.  The same initial-segment chain now puts $A$ and every
earlier hierarchy stage inside $T(y_{\ell-1})$ entry by entry; restricting
\noderef{ColoredVertexSetElementary} gives the claim.  In particular,
whenever a color in the proof is missing on one of these stages, its
first missing vertex there is unique.

Second, swapping a full component preserves properness by
\noderef{TwoColorComponentSwapProper}, applied to the given properness
and $\alpha\ne\beta$.  Unfolding
\noderef{TwoColorComponentSwap} shows separately that the value
$\mathsf{none}$ is fixed and only the two nonempty displayed colors are
transposed.  Therefore an edge is uncolored under $\sigma_n$ exactly
when it is uncolored under $\varphi_n$; the given
\noderef{IsOneEdgeUncolored} property shows that precisely $e_0$
remains uncolored.  We shall use these two properties of $\sigma_n$
without repeating the argument.

Put $a=v_\alpha$ and $b=v_\beta$.  Since every vertex of $A$ precedes
$a$, the least-candidate definition in
\noderef{HierarchyFirstMissingPosition} gives
\[
             \alpha,\beta\notin\overline\varphi_n(A).         \tag{6}
\]
Because (6) says that neither color is missing on $A$, the definition
\noderef{HierarchyStoredColorsAt} turns
$\alpha,\beta\notin D_{n,q}$ into
$\alpha,\beta\notin D_n$.  The statement of
\noderef{HkuTreeColorsControlled} puts every color on a displayed edge
of $T_n$ in $\overline\varphi_n(T_n)\cup D_n$; (6) excludes the first
set and the preceding observation excludes the second.  Thus neither
color occurs on $T_n$.  Every later ordinary entry before $b$ is a
\noderef{TashkinovAugmentation} edge, so an $\alpha$- or
$\beta$-colored entry would have an earlier missing witness, contrary to
the definitions of $a$ and $b$.  A possible SE/RE connecting seed has
its color in $D_n$ by \noderef{TashkinovSeriesStep}, so it is excluded
as well.  Combining these facts with (U) gives the identity exclusions
\[
 \alpha,\beta\notin\varphi_n\langle T(v_\beta)\rangle,
 \qquad
 \alpha,\beta\notin
   \overline\varphi_n(A)\cup D_{n,q}.                         \tag{7}
\]
No endpoint pair has replaced the indexed edges in (U).

The hypotheses of \noderef{ChangeLemma} are now all available.  Its
ambient lower-rung, fewer-divider, same-rung minimality, prefix-elementarity, and
large-color assumptions are exactly those stated here.  The order
$a\prec b$ is given.  If one of the colors is not missing on $A$, (6)
says both are not missing there.  Any $\alpha$-colored ordinary
augmentation edge between $A$ and $a$ would, by
\noderef{TashkinovAugmentation}, have an earlier $\alpha$-missing
witness and contradict the minimality of $a$; the possible stored seed
is excluded by $\alpha\notin D_{n,q}$.  Combining this with (U) gives
the required exclusion from $A$ to $b$.  Finally, if $q>0$, (7) gives the
disjoint-$D_{n,q}$ alternative of its case 1; if $q=0$, (6) and the
definition of \noderef{HierarchyStoredColorsAt} turn (7) into
$\alpha,\beta\notin D_n$, giving its case 2.  Consequently
\[
 K_a=K_b                                                       \tag{8}
\]
as full incidence components.

We verify stability without assuming that $Q$ is different from the
component in (8), and without assuming that an unrelated component is
outside $A$.  By (6), every vertex of $A$ has both $\alpha$ and $\beta$
present.  Properness gives exactly one incident edge of each color
there.  If $Q$ meets such a vertex, both incident edge-nodes belong to
$Q$ and their colors are exchanged together; hence its missing-color
set is unchanged.  If $Q$ does not meet the vertex, nothing changes.
Moreover (6)--(7) and the preceding $D_n$ calculation show that neither
affected color lies in the controlling set
$\overline\varphi_n(T_n)\cup D_n$, neither is missing on $R_n$, and no
displayed edge of $A$ has an affected color.  Unfolding
\noderef{StableColoring} therefore gives simultaneous stability on
$T_n$ and, in the PE case, on $R_n$, while the displayed-edge and vertex
clauses give invariance on $A$.  This is precisely
\noderef{HierarchyWeaklyStable} at the last starred stage.

We next extract the terminal construction with all of its exact seed
data.  Let $\widehat T$ be the last closure tree in the certified series.
The positive terminal instance of \noderef{TashkinovSeriesStep} has
exactly one of the following forms.
\begin{itemize}
\item If $\Theta_n=PE$, its closure field is a
$\varphi_n$-\noderef{TashkinovClosure} from the exact ordered tree $T_n$
to $\widehat T$.
\item If $\Theta_n=SE$ or $RE$, it supplies an individual edge identity
$f_n$, the exact equality $F_n=\{f_n\}$, an exact
\noderef{TreeSequenceAdjoinEdge} seed $S=T_n+f_n$, and a
$\varphi_n$-closure from $S$ to $\widehat T$.
\end{itemize}
In either case, unfolding \noderef{TashkinovClosure} gives the displayed
augmentation sequence from its stated seed to $\widehat T$.  Each
one-edge step has its predecessor as an exact ordered
\noderef{TreeSequenceInitialSegment} by
\noderef{TashkinovAugmentation}; composing the entrywise vertex and edge
equalities makes the seed an exact ordered initial segment of
$\widehat T$.

The ETT certificate makes $T$ an exact ordered initial segment of
$\widehat T$.  The base-stage, nested-stage, and terminal-stage fields of
\noderef{TashkinovHierarchy} make $A$ an exact ordered initial segment of
$T$, hence of $\widehat T$.  At $q=0$ this uses the two alternatives in
\noderef{HierarchyBase}: in PE, $A=T_n\vee R_n$ and the left part of
\noderef{OrderedTreeUnion} is the exact prefix $T_n$; in SE/RE,
$A=T_n$.  At $q>0$, $A=H.\mathit{stage}(q)$ and the same conclusion is
the composite of the hierarchy's exact nested-stage maps.

It follows that the terminal seed is an exact prefix of every target to
which it must be restricted.  In PE the seed is $T_n$, already an exact
prefix of both $A$ and $T$.  In SE/RE, the seed $S$ has length
$|T_n|+1$.  If $q>0$, strictness of the first hierarchy stage gives
$|T_n|+1\leq |H.\mathit{stage}(1)|\leq |A|$.  Both $S$ and $A$ are exact
prefixes of $\widehat T$, so their entries through $|S|$ are the same
entries of $\widehat T$; hence $S$ is an exact ordered prefix of $A$.
The identical comparison with $T$, using the strict ETT inequality
$|T_n|<|T|$, makes $S$ an exact prefix of $T$.  Thus
\noderef{TashkinovAugmentationSequenceRestrict}, first with intermediate
tree $A$ when $A\ne T_n$ and then with intermediate tree $T$, gives the
exact old-coloring routes
\[
 \begin{array}{ll}
 T_n\longrightarrow A,\ T_n\longrightarrow T,&\Theta_n=PE,\\
 T_n+f_n\longrightarrow A,\ T_n+f_n\longrightarrow T,
     &\Theta_n\in\{SE,RE\},\ q>0,
 \end{array}                                                    \tag{9}
\]
and the second SE/RE route also holds when $q=0$.  When $q=0$ in the
non-PE case there is intentionally no route to repackage $A$, because
$A=T_n$ is the endpoint alternative in conclusion 2.

We now transport every one-edge step in (9) under $\sigma_n$.  Any such
route adds the displayed identities in their order in $T$.  Through the
prefix $T(b)$, equation (7) says that no new or previously displayed
identity has old color $\alpha$ or $\beta$.  The component swap therefore
fixes its color.  A color $c\notin\{\alpha,\beta\}$ is missing at a
vertex under $\sigma_n$ exactly when it is missing there under
$\varphi_n$, since the swap changes only the two affected incident
colors.  Hence every augmentation through $T(b)$ retains the same new
edge color and the same missing-color witness.

Consider a later augmentation, whose predecessor already contains both
$a$ and $b$.  The definitions of the first-missing positions give that
$\alpha$ is missing at $a$ and $\beta$ is missing at $b$.  Since both
vertices lie in the elementary prefix $T(y_{\ell-1})$, $\beta$ is present at
$a$ and $\alpha$ is present at $b$; properness makes the relevant present
edge unique.  By (8) their common component is $K=K_a=K_b$.
\noderef{TwoColorComponentsDisjoint} says that the swapped component $Q$
is either $K$ or is disjoint from it.

If $Q=K$, the swap makes $\beta$ missing at $a$ and $\alpha$ missing at
$b$.  If $Q\ne K$, neither $a$ nor $b$, nor any affected edge incident
with either vertex, belongs to $Q$, so $\alpha$ remains missing at $a$
and $\beta$ remains missing at $b$.  Now let $f$ be the exact new edge of
the augmentation.  If its old color is outside
$\{\alpha,\beta\}$, both its color and its old missing witness are fixed.
If its old color is $\alpha$ or $\beta$, inspect whether the edge-node
$f$ lies in $Q$.  Inside $Q$ its color is transposed and outside $Q$ it
is fixed.  In the case $Q=K$, choose respectively the post-swap witness
among $a,b$ just recorded; in the case $Q\ne K$, choose $a$ for color
$\alpha$ and $b$ for color $\beta$, exchanging these choices when $f$
lies in $Q$ and its color is transposed.  In all four combinations the
post-swap color of $f$ is missing at the chosen one of $a,b$, and that
vertex precedes $f$.  The exact length, attachment, and ordered-prefix
clauses are independent of the coloring.  This verifies every clause of
\noderef{TashkinovAugmentation} at every retained step, and hence turns
all applicable routes in (9) into
$\sigma_n$-\noderef{TashkinovAugmentationSequence}s.             \tag{10}

The weak stability already proved contains its base-stability clause.
In the PE case that clause is precisely
\noderef{CombinedStableColoring} on $T_n$ and $R_n$; in the SE/RE case
it is $(T_n,D_n,\varphi_n)$-stability.  In particular, in all cases it
contains the latter terminal-ladder stability.  Apply \noderef{Td} with
the stated six lower-rung assertions, the original maximum property, and
the properness and unique-uncolored facts proved above.  We obtain
\[
                      \noderef{ModColoring}(\mathit{ett},\sigma_n). \tag{11}
\]
This proves the first assertion in conclusion 2.  If $A=T_n$, its first
alternative holds.  Otherwise $T_n\prec A$, and (10) supplies exactly the
PE route from $T_n$, or the exact SE/RE seed $T_n+f_n$ and route, required
by the universal clause of (11).  That clause gives an actual
\noderef{CorrespondingExtendedTashkinovTree} on $A$, and
\noderef{MaximumPropertyTransport} gives its maximum property.  This
proves the rest of conclusion 2.

We next construct the full-tree hierarchy certificate under $\sigma_n$.
First suppose the terminal kind is PE.  Apply
\noderef{HierarchySplitter} to the original maximum ETT and hierarchy.
Its combined-stability premise is the PE projection of the already proved
\noderef{HierarchyWeaklyStable}; its properness and one-uncolored-edge
premises were proved above; and its final augmentation premise is the
$T_n\longrightarrow T$ route in (10).  The lower-rung package is exactly
the one stated here.  The lemma returns an actual corresponding ETT on
$T$, maximum property, a hierarchy base under $\sigma_n$, and a hierarchy
$K$ with the same exact starred core, closure, ordered stages, and divider
positions as $H$.

If the terminal kind is SE or RE, use the exact seed and route in (10) in
the universal clause of (11).  It returns an actual corresponding ETT on
$T$, and \noderef{MaximumPropertyTransport} again gives maximum property.
Rebuild its \noderef{HierarchyBase} field by field: correspondence gives
the same terminal ladder tree and the same non-PE kind; take the core to
be the common exact tree $T_n$, copy the unused closure and the positive
support position, use reflexivity for the terminal missing-color equation
under the generating coloring $\sigma_n$, and use the common non-PE kind
for the nonparallel field.  Copy $H$'s divider count, every exact ordered
stage, and every divider position.  The base-stage equation now has the
same two sides $T_n$, while terminality, nestedness, strictness, divider
prefixes, and core-inside-first are literally the old fields.  This gives
a hierarchy $K$ with the same stages and dividers as $H$.  In both
branches, denote the resulting corresponding ETT, maximum witness, base,
and hierarchy by $\mathcal E^\sigma,m^\sigma,B^\sigma,K$; their base core
and closure equal those of the original base, including the PE closure
equality returned by \noderef{HierarchySplitter}.               \tag{12}

It remains to construct, rather than assert, the good-hierarchy witness
on $K$.  Use the equality of divider counts in (12) to identify its stage
indices with those of $H$, and define
\[
                    \Gamma^{\sigma,j}_\eta=\Gamma^j_\eta       \tag{13}
\]
for every value of the total Gamma function.  Every starred stage and
every ordinary stage occurring at an index $j\leq q$ is an exact ordered
subtree of $A$.  At each of its vertices both affected colors are present
by (6).  If the vertex-node lies in $Q$, properness puts its unique
$\alpha$- and $\beta$-edges in the same component and the swap exchanges
them together; if it does not lie in $Q$, neither incident affected edge
is changed.  All other colors are fixed.  Therefore the complete missing
set at every such vertex is the same under $\varphi_n$ and $\sigma_n$.
Taking unions proves equality of all the controlled sets in
\noderef{HierarchyStageMissingColors} and of the missing-color sets of
all starred stages.                                                   \tag{14}

The corresponding ETT certificate identifies every raw set $S_h$ through
the terminal index, while (14) identifies the terminal-ladder missing
sets.  The union-and-difference definition of
\noderef{TashkinovStoredColors} therefore gives the same terminal stored
set; (14) also identifies the starred-stage set subtracted from it in
\noderef{HierarchyStoredColorsAt}.  Consequently
\[
                  D^K_{n,j}=D^H_{n,j}\quad(0\leq j\leq q).     \tag{15}
\]
The earlier deduction from (6)--(7) gave
$\alpha,\beta\notin D_n$, while every set in (15) is a subset of $D_n$.
Thus neither affected color is an index in any set (15).  More
importantly, the avoidance field of the old \noderef{GoodHierarchy} and
(14) say that every color in every copied pair (13) is missing on a
controlled stage.  Since neither $\alpha$ nor $\beta$ is missing on such
a stage, neither affected color belongs to any copied Gamma pair.       \tag{16}

We check the five fields of \noderef{GoodHierarchy}.
Pair cardinality and pairwise disjointness follow immediately from
(13)--(15).  For avoidance, fix $j,\eta$ in the common stored set and
$c\in\Gamma^j_\eta$.  By (16), $c\notin\{\alpha,\beta\}$.  The indexed
color also satisfies $\eta\notin\{\alpha,\beta\}$, by (15) and
$\alpha,\beta\notin D_n$.  Swapping the two affected colors preserves,
vertex by vertex on the full tree, whether $\eta$ is missing.  Hence the
candidate set and the totalized value of
\noderef{HierarchyFirstMissingPosition} for $\eta$ are unchanged.  The old and
new $(j,\eta)$ segments in \noderef{HierarchySegmentColors} therefore
contain exactly the same individual edge positions.  If one of those
edges had new color $c$, then, since $c$ is unaffected by the swap, it
would have old color $c$, contrary to the old avoidance field.  This is
the required pointwise avoidance proof.  It does not assert equality of
the full segment-color sets: their irrelevant $\alpha$- and $\beta$-parts
may indeed be exchanged.

For unique refresh, the two relevant ordered stages and starred stage are
the same by (12), their vertex-set difference is the same, and (14)
identifies the missing-color set at each vertex in that difference.
Together with (13)--(15), the old existence-and-uniqueness witness is
therefore exactly a witness for the new field.

For the parallel-base field, suppose the common terminal kind is PE.
The base-stability clause is \noderef{CombinedStableColoring}; its two
vertex clauses identify the missing sets on $T_n$ and on $R_n$.
Equations (13) and (15) identify the reserved union $\Gamma^0$, so the
old common-missing cardinality inequality is unchanged.  To transport
non-minus-closedness, unpack the old witness: there are a color $c$
missing on the old core but outside $\Gamma^0$ and an individual boundary
edge $g$ of that core with $\varphi_n(g)=c$.  The core lies in $A$, so
(6) forces $c\notin\{\alpha,\beta\}$.  Thus the swap fixes $g$ and gives
$\sigma_n(g)=c$; (14) preserves the missingness of $c$, and (12)
preserves the exact core and its boundary.  The same $c,g$ witness that
the new core is not $(\Gamma^0)^-$-closed.  This proves both PE clauses.
In the common non-PE case the implication defining the parallel-base
field has a false PE premise and is immediate.

Finally fix an internal index $i<q$.  Put
\[
 U=H.\mathit{stage}(i+1)=K.\mathit{stage}(i+1),\qquad
 R=\bigcup_{\eta\in D_{n,i+1}}\Gamma^i_\eta .
\]
To verify the new minus-closed field, let $c$ be missing on $U$ under
$\sigma_n$ with $c\notin R$, and suppose that an individual boundary edge
$g$ of $U$ has $\sigma_n(g)=c$.  Since $U\subseteq A$, neither affected
color is missing there, so $c\notin\{\alpha,\beta\}$.  Equation (14) says
$c$ is missing on $U$ under $\varphi_n$, (13)--(15) give the same reserved
set $R$, and the swap rule gives $\varphi_n(g)=c$.  This contradicts the
old minus-closed field.  Thus every internal minus-closed field holds for
$K$.  Notice that its final instance, at $i=q-1$, has been proved here
once and will be reused below.  We have now checked every field and obtain
a good hierarchy $\Gamma^\sigma$ on $K$.  Equations (12), (15), and (13)
are exactly \noderef{SameGoodHierarchy}$(C.\mathit{good},\Gamma^\sigma)$,
which proves the full-tree conclusion.                                \tag{17}

Assume at last that $q>0$.  Apply
\noderef{GoodHierarchyPrefixRepackage} to
$\mathcal E^\sigma,m^\sigma,B^\sigma,K,\Gamma^\sigma$ at its last
positive internal stage $r=q$.  Because $q>0$, the definition of
\noderef{HierarchyStarredStage}, followed by the exact stage equality in
(12), identifies
\[
 K.\mathit{stage}(r)=H.\mathit{stage}(q)=A=T_{n,q}.
\]
The theorem therefore returns an actual maximum ETT on the exact ordered
tree $A$, a base with the same core and closure, and a good hierarchy with
divider count $q-1$ and exactly the retained stages through $A$.  Its
retained stored-set and total Gamma equalities, composed with (13)--(15),
give precisely the $D_{n,j}$ and paper Gamma equalities required in
conclusion 3.

The prefix theorem also returns separately the final internal
minus-closed field of the full hierarchy (17).  Under the same exact stage,
stored-set, and Gamma equalities, that returned assertion is
\[
 T_{n,q}\text{ is }
 \left(\bigcup_{\eta\in D_{n,q}}\Gamma^{q-1}_\eta\right)^-
 \text{-closed under }\sigma_n,
\]
which is (M).  Thus (M) is the already verified last boundary argument,
not a second unsupported transport claim.  All three conclusions and the
full-tree certificate now follow.
\end{proof}
